% !TeX root = ../../../../../main.tex
% chapters/sections/chapter3/subsections/resource_constraints/aggregate.tex

\subsection{自国の資源制約}
\label{sec:chapter3_resource_constraints_aggregate_home}
式
\eqref{form:chapter3_households_stage2_home_budget_eq}
で示されたように、自国家計\(h\)の予算制約式は以下のようであった。
\begin{equation}
\label{form:chapter3_resource_constraints_aggregate_home_budget_eq}
    \sum_{j'\in J}q_{t,t+1}(j')d_{t+1}^{h\to H}(j')+b_{t+1}^{h\to F}+p_t^{H\to W}c_t^{h\to W}
    =d_t^{h\to H}+(1+i_{t-1}^F)\frac{e_t^{/*}}{e_{t-1}^{/*}}b_t^{h\to F}+(1-\tau_t^H)p_t^hy_t^h+t_t^H
\end{equation}
この予算制約式をすべての自国家計\(h\)について足し合わせることで自国全体の資源制約式を導出する。
まずすべての自国家計\(h\)について予算制約式の両辺を足し合わせると以下の式が得られる。
\begin{equation}
\label{form:chapter3_resource_constraints_aggregate_home_sum_budgets_eq}
    \sum_{h\in H}\left(\sum_{j'\in J}q_{t,t+1}(j')d_{t+1}^{h\to H}(j')+b_{t+1}^{h\to F}+p_t^{H\to W}c_t^{h\to W}\right)
    =\sum_{h\in H}\left(d_t^{h\to H}+(1+i_{t-1}^F)\frac{e_t^{/*}}{e_{t-1}^{/*}}b_t^{h\to F}+(1-\tau_t^H)p_t^hy_t^h+t_t^H\right)
\end{equation}
さらにこの和の記号を各項に分配し、家計\(h\)に依存しない変数（物価水準、利子率、為替レート、税率など）を和の記号の外にくくり出す。
その際、左辺第1項のコンティンジェント債券の購入費用については、状態価格\(q_{t,t+1}(j')\)が家計に依存しないため、和の順序を入れ替えてくくり出すことができる。整理すると以下の展開式が得られる。
\begin{align}
\label{form:chapter3_resource_constraints_aggregate_home_sum_budgets_modified_eq}
    &\sum_{j'\in J}q_{t,t+1}(j')\sum_{h\in H}d_{t+1}^{h\to H}(j')+\sum_{h\in H}b_{t+1}^{h\to F}+p_t^{H\to W}\sum_{h\in H}c_t^{h\to W} \nonumber \\
    &\quad =\sum_{h\in H}d_t^{h\to H}+(1+i_{t-1}^F)\frac{e_t^{/*}}{e_{t-1}^{/*}}\sum_{h\in H}b_t^{h\to F}+(1-\tau_t^H)\sum_{h\in H}p_t^hy_t^h+\sum_{h\in H}t_t^H
\end{align}
次に、この展開式の各項に代入して整理するための条件を準備する。
まず式
\eqref{form:chapter3_assumptions_bond_markets_home_equilibrium}
の自国コンティンジェント債券市場の均衡条件より、
次期および今期の自国コンティンジェント債券の総和はそれぞれゼロとなる。
\begin{equation}
\label{form:chapter3_resource_constraints_aggregate_home_equilibrium_eq}
    \sum_{h\in H}d_{t+1}^{h\to H}(j')=0\quad (\forall j'\in J)
\end{equation}
\begin{equation}
\label{form:chapter3_resource_constraints_aggregate_home_equilibrium2_eq}
    \sum_{h\in H}d_t^{h\to H}=0
\end{equation}
次に式
\eqref{form:chapter3_policies_fiscal_home_budget_eq}
の政府予算制約より、一括移転の総和\(\sum_{h\in H}t_t^H\)は以下のようになる。
\begin{equation}
\label{form:chapter3_resource_constraints_aggregate_home_fiscal_budget_eq}
    \sum_{h\in H}t_t^H=Nt_t^H=\tau_t^H\sum_{h\in H}p_t^hy_t^h
\end{equation}
さらに式
\eqref{form:chapter3_producers_aggregated_production_home_Y_H_definition_p_H_Y_H_eq}
で示されたように、名目総生産\(\sum_{h\in H}p_t^hy_t^h\)は以下のように書き換えられる。
\begin{equation}
\label{form:chapter3_resource_constraints_aggregate_home_p_H_Y_H_eq}
    \sum_{h\in H}p_t^hy_t^h=p_t^HY_t^H
\end{equation}
また自国の外国通貨建てリスクフリー債券\(B_t^H\)および総消費\(C_t^{H\to W}\)を以下のように定義する。
\begin{equation}
\label{form:chapter3_resource_constraints_aggregate_home_B_H_def}
    B_t^H\equiv\sum_{h\in H}b_t^{h\to F}
\end{equation}
\begin{equation}
\label{form:chapter3_resource_constraints_aggregate_home_C_H_W_def}
    C_t^{H\to W}\equiv\sum_{h\in H}c_t^{h\to W}
\end{equation}
これらを代入すると以下の資源制約式が得られる。
\begin{align}
    \label{form:chapter3_resource_constraints_aggregate_home_sum_budgets_modified2_eq}
        0+B_{t+1}^H+p_t^{H\to W}C_t^{H\to W}
        &=0+(1+i_{t-1}^F)\frac{e_t^{/*}}{e_{t-1}^{/*}}B_t^H+(1-\tau_t^H)p_t^HY_t^H+\tau_t^Hp_t^HY_t^H \\
    \label{form:chapter3_resource_constraints_aggregate_home_sum_budgets_modified3_eq}
        B_{t+1}^H+p_t^{H\to W}C_t^{H\to W}
        &=p_t^HY_t^H+(1+i_{t-1}^F)\frac{e_t^{/*}}{e_{t-1}^{/*}}B_t^H
\end{align}
この式は自国の外国通貨建てリスクフリー債券と総消費が
名目生産と前期からの外国通貨建てリスクフリー債券の収益によって賄われることを示している。

\subsection{外国の資源制約}
\label{sec:chapter3_resource_constraints_aggregate_foreign}
式
\eqref{form:chapter3_households_stage2_foreign_budget_eq}
で示されたように、外国家計\(f\)の予算制約式は以下のようであった。
\begin{equation}
\label{form:chapter3_resource_constraints_aggregate_foreign_budget_eq}
    \sum_{j'\in J}q_{t,t+1}^*(j')d_{t+1}^{f\to F*}(j')+b_{t+1}^{f\to H*}+p_t^{F\to W*}c_t^{f\to W}
    =d_t^{f\to F*}+(1+i_{t-1}^H)\frac{e_{t-1}^{/*}}{e_t^{/*}}b_t^{f\to H*}+(1-\tau_t^F)p_t^{f*}y_t^f+t_t^{F*}
\end{equation}
この予算制約式をすべての外国家計\(f\)について足し合わせることで外国全体の資源制約式を導出する。
まずすべての外国家計\(f\)について予算制約式の両辺を足し合わせると以下の式が得られる。
\begin{equation}
\label{form:chapter3_resource_constraints_aggregate_foreign_sum_budgets_eq}
    \sum_{f\in F}\left(\sum_{j'\in J}q_{t,t+1}^*(j')d_{t+1}^{f\to F*}(j')+b_{t+1}^{f\to H*}+p_t^{F\to W*}c_t^{f\to W}\right)
    =\sum_{f\in F}\left(d_t^{f\to F*}+(1+i_{t-1}^H)\frac{e_{t-1}^{/*}}{e_t^{/*}}b_t^{f\to H*}+(1-\tau_t^F)p_t^{f*}y_t^f+t_t^{F*}\right)
\end{equation}
さらにこの和の記号を各項に分配し、家計\(f\)に依存しない変数を和の記号の外にくくり出す。
状態価格\(q_{t,t+1}^*(j')\)の和の順序を入れ替えて整理すると、以下の展開式が得られる。
\begin{align}
\label{form:chapter3_resource_constraints_aggregate_foreign_sum_budgets_modified_eq}
    &\sum_{j'\in J}q_{t,t+1}^*(j')\sum_{f\in F}d_{t+1}^{f\to F*}(j')+\sum_{f\in F}b_{t+1}^{f\to H*}+p_t^{F\to W*}\sum_{f\in F}c_t^{f\to W} \nonumber \\
    &\quad =\sum_{f\in F}d_t^{f\to F*}+(1+i_{t-1}^H)\frac{e_{t-1}^{/*}}{e_t^{/*}}\sum_{f\in F}b_t^{f\to H*}+(1-\tau_t^F)\sum_{f\in F}p_t^{f*}y_t^f+\sum_{f\in F}t_t^{F*}
\end{align}
次に、この展開式の各項に代入して整理するための条件を準備する。
まず式
\eqref{form:chapter3_assumptions_bond_markets_foreign_equilibrium}
の外国コンティンジェント債券市場の均衡条件より、
次期および今期の外国コンティンジェント債券の総和はそれぞれゼロとなる。
\begin{equation}
\label{form:chapter3_resource_constraints_aggregate_foreign_equilibrium_eq}
    \sum_{f\in F}d_{t+1}^{f\to F*}(j')=0\quad (\forall j'\in J)
\end{equation}
\begin{equation}
\label{form:chapter3_resource_constraints_aggregate_foreign_equilibrium2_eq}
    \sum_{f\in F}d_t^{f\to F*}=0
\end{equation}
次に式
\eqref{form:chapter3_policies_fiscal_transfer_foreign_eq}
の政府予算制約より、一括移転の総和\(\sum_{f\in F}t_t^{F*}\)は以下のようになる。
\begin{equation}
\label{form:chapter3_resource_constraints_aggregate_foreign_fiscal_budget_eq}
    \sum_{f\in F}t_t^{F*}=\tau_t^F\sum_{f\in F}p_t^{f*}y_t^f
\end{equation}
さらに式
\eqref{form:chapter3_producers_aggregated_production_foreign_Y_F_definition_p_F_Y_F_eq}
で示されたように、名目総生産\(\sum_{f\in F}p_t^{f*}y_t^f\)は以下のように書き換えられる。
\begin{equation}
\label{form:chapter3_resource_constraints_aggregate_foreign_p_F_Y_F_eq}
    \sum_{f\in F}p_t^{f*}y_t^f=p_t^{F*}Y_t^F
\end{equation}
また外国の自国通貨建てリスクフリー債券\(B_t^{F*}\)および総消費\(C_t^{F\to W}\)を以下のように定義する。
\begin{equation}
\label{form:chapter3_resource_constraints_aggregate_foreign_B_F_def}
    B_t^{F*}\equiv\sum_{f\in F}b_t^{f\to H*}
\end{equation}
\begin{equation}
\label{form:chapter3_resource_constraints_aggregate_foreign_C_F_W_def}
    C_t^{F\to W}\equiv\sum_{f\in F}c_t^{f\to W}
\end{equation}
これらを代入すると以下の資源制約式が得られる。
\begin{align}
    \label{form:chapter3_resource_constraints_aggregate_foreign_sum_budgets_modified2_eq}
        0+B_{t+1}^{F*}+p_t^{F\to W*}C_t^{F\to W}
        &=0+(1+i_{t-1}^H)\frac{e_{t-1}^{/*}}{e_t^{/*}}B_t^{F*}+(1-\tau_t^F)p_t^{F*}Y_t^F+\tau_t^Fp_t^{F*}Y_t^F \\
    \label{form:chapter3_resource_constraints_aggregate_foreign_sum_budgets_modified3_eq}
        B_{t+1}^{F*}+p_t^{F\to W*}C_t^{F\to W}
        &=p_t^{F*}Y_t^F+(1+i_{t-1}^H)\frac{e_{t-1}^{/*}}{e_t^{/*}}B_t^{F*}
\end{align}
この式は外国の自国通貨建てリスクフリー債券と総消費が
名目生産と前期からの自国通貨建てリスクフリー債券の収益によって賄われることを示している。